% Part: proof-theory
% Chapter: natural-deduction
% Section: introduction

\documentclass[../../../include/open-logic-section]{subfiles}

\begin{document}

\olfileid[tr]{pt}{nat}{int}

\olsection{Giriş}

Doğal tümdengelim, her mantıksal bağlaç ya da niceleyici için biri işleci
``tanıtan'' (!!a{introduction} kuralı), öteki işleci ``gideren''
(!!a{elimination} kuralı) bir çıkarım çifti bulunan !!{proof} sistemleri
ailesidir. Örneğin \Intro{\lif}, sonuç olarak~$!A \lif !B$ biçimindeki
!!{formula}{}ü; \Elim{\lif} ise öncül olarak~$!A \lif !B$ biçimindeki
!!{formula}{}ü içerir.

Doğal tümdengelimin ilk iki sürümü 1930'larda Gerhard Gentzen ve Stanisław
Jaśkowski tarafından tanıtıldı. İspat kuramcıları daha çok Gentzen'in
sistemini ele alır; öğretimde en yaygın kullanılan sistemler ise
Jaśkowski'nin yaklaşımından doğmuştur. İki sistemde de \emph{varsayımlarla}
başlanabilir: Bunlar !!{prove}{}nması gerekmeyen, fakat başka !!{formula}s{}i
!!{proving} amacıyla kullanılabilen !!{formula}s{}dir. Bütün bir !!{proof} böyle
varsayımlara dayanabilir. Varsayımlar ispatlanmadığından !!{proof}, sonucunu
(\emph{son-!!{formula}{}ünü}) tek başına değil, yalnızca sonucun varsayımlardan
çıktığını gösterir.

Bazı çıkarım kuralları, sonuçlarının artık kimi varsayımlara dayanmamasını
sağlar: Varsayımları ``kapatırlar.'' Örneğin \Intro{\lif} kuralı, $!A$
varsayımından $!B$ sonucuna giden !!a{proof}{}ı alır ve sonuç olarak
$!A \lif !B$'yi verir. $!A \lif !B$ ile biten !!{proof} artık $!A$'ya
dayanmaz; $!A$ varsayımı kapatılmıştır. Gentzen sistemlerinde !!{proof}s,
!!{formula}s{}den oluşan ağaçlardır; varsayımlar en üstteki !!{formula}s ve
\Intro{\lif} gibi kurallar hangi varsayımların kapatıldığını belirten
etiketler taşır. Jaśkowski sistemlerinde bir varsayıma dayanan !!{proof}
parçaları bir biçimde ayrılır: Örneğin dikey bir çizginin gerisine girintili
yazılır ya da bir kutuya alınır. Kapatılacak $!A$ varsayımı kutunun ilk,
koşullunun artbileşeni $!B$ son satırında; \Intro{\lif}'in sonucu
$!A \lif !B$ ise kutunun dışında bulunur.

Bu sistemler ``doğal''dır; çünkü çıkarım kuralları bizim (en azından
matematikçilerin) doğal akıl yürütmesini yansıtır. Varsayımsal bir sonucu
(koşulluyu) göstermek için önbileşeni varsayar ve artbileşeni ispatlamakta
kullanırız. Bu~\Intro{\lif}'dir. Bir $!A \lif !B$ koşullusuyla önbileşeni
$!A$ bilindiğinde $!B$ sonucuna varırız. Bu~\Elim{\lif}'dir. Bir ayrışım
$!A \lor !B$ varsayımı altında bir sonucun geçerli olduğunu göstermek için
durumlara ayırarak ispat kullanırız. Bu~\Elim{\lor}'dir. Genel bir
$\lforall[x][!A(x)]$ iddiasını ispatlamak için gelişigüzel bir~$c$ için
$!A(c)$'nin geçerli olduğunu ispatlarız. Bu~\Intro{\lforall}'dır. Ve böyle
devam eder.

Örneğin Gentzen tipi doğal tümdengelimde $!A \lif (!A \lor !B)$'nin basit
bir !!{proof}{}ı şöyledir:
\[
\AxiomC{$\Discharge{!A}{x}$}
\RightLabel{\Intro{\lor}}
\UnaryInfC{$!A \lor !B$}
\DischargeRule{\Intro{\lif}}{x}
\UnaryInfC{$!A \lif (!A \lor !B)$}
\DisplayProof
\]
$!A$ varsayımıyla başlar,~$!A \lor !B$ sonucunu çıkarmak için
\Intro{\lor}, ardından~$!A \lif (!A \lor !B)$ sonucunu çıkarmak için
\Intro{\lif} kullanılır. \Intro{\lif} çıkarımı $!A$ varsayımını kapatır;
bunu~$x$ etiketi gösterir.

Bir çeşitlemesi de ispat kuramında önemli olsa da ispat kuramcıları
Gentzen'in !!{formula}s ağaçlarıyla çalışan özgün sistemlerini yeğler.
$\Gamma$ varsayımlarından~$!A$'nın !!^a{proof}{}ı, $!A$'nın $\Gamma$'dan
çıktığını, yani $\Gamma \Entails !A$ olduğunu gösterir. Doğal tümdengelimin
çıkarım kuralları böyle gerektirmeler hakkında bir şey söylediği biçiminde
doğal olarak okunabilir. Örneğin \Intro{\lif},
$\Gamma + !A \Entails !B$ ise $\Gamma \Entails !A \lif !B$ olduğunu söyler.
Kuralları, çıkarım kuralının öncülleriyle sonucunda hem varsayımlar hem de
onlardan çıkan sonuç üzerinde, yani $\Gamma \Sequent !A$ sekantları üzerinde
doğrudan işleyecek biçimde formüle etmek doğaldır. Yukarıdaki basit ispat
böyle bir sistemde şöyle görünür:
\[
\Axiom$x:!A \fCenter !A$
\RightLabel{\Intro{\lor}}
\UnaryInf$x:!A \fCenter !A \lor !B$
\DischargeRule{\Intro{\lif}}{x}
\UnaryInf$\fCenter !A \lif (!A \lor !B)$
\DisplayProof
\]
Gördüğünüz gibi kurallar aynıdır: $\Sequent$ simgesinin sağındaki
!!{formula} üzerinde işler; soldaki !!{formula}s ise yalnızca her adımın
dayandığı varsayımları listeler.

!!{introduction}/!!{elimination} kural çiftlerinin tam kümesi tek başına
klasik mantık için tam değildir. $\lfalse$ ile ilgili iki ek kural
eklemeliyiz. Birincisi, $\lfalse$'dan her şeyin çıkarılabileceğini söyleyen
``sezgici saçmalık kuralı''~\FalseInt, diğer adıyla ``patlama''dır. İkincisi,
$\lnot !A$'dan~$\lfalse$'ı !!{prove}{}yabilirsek $!A$'yı !!{prove}{}mış
olduğumuzu söyleyen klasik dolaylı ispat ilkesi~\FalseCl{} (ya da ``klasik
saçmalık kuralı'')dır. \FalseCl{} çıkarılırsa sezgici mantığa uygun bir sistem
elde ederiz. İkisi de çıkarılırsa ``minimal mantık'' denen sistem kalır.

Klasik ve sezgici mantık için Gentzen tipi doğal tümdengelimi ele alacağız.
Bunlar $\Log{N1c}$ ve $\Log{N1i}$ sistemleridir. $\Log{N2c}$ ile
$\Log{N2i}$ ise tek tek !!{formula}s~$!A$ yerine $\Gamma \Sequent !A$
sekantları üzerinde çalışan karşılık gelen sistemlerdir.

\end{document}
